%----------
%	REGULATORY FRAMEWORK
%----------	

En esta sección se detallan las bases normativas y los principios éticos que enmarcan el desarrollo de la presente investigación.

En primer lugar, dado que el análisis de textos periodísticos conlleva el procesamiento de información que puede identificar a personas físicas (firmas de periodistas, nombres de fuentes expertas o figuras públicas), el proyecto se supedita estrictamente a las directrices del Reglamento General de Protección de Datos (RGPD)\footnote{Reglamento (UE) 2016/679 del Parlamento Europeo y del Consejo. Disponible en: \url{https://eur-lex.europa.eu/legal-content/ES/TXT/?uri=CELEX\%3A32016R0679} [Último acceso: 15/03/2026]}. La metodología garantiza que el tratamiento de estos datos se realiza con un fin analítico centrado en detectar patrones lingüísticos y sesgos de género, eludiendo cualquier tipo de perfilado personal de la información extraída.

En el ámbito tecnológico, el uso de los LLM y de sistemas basados en agentes obliga a contextualizar este TFM dentro del Reglamento de Inteligencia Artificial de la Unión Europea (AI Act)\footnote{Reglamento (UE) 2024/1689 del Parlamento Europeo y del Consejo, de 13 de junio de 2024, por el que se establecen normas armonizadas en materia de inteligencia artificial. Disponible en: \url{https://eur-lex.europa.eu/legal-content/ES/TXT/?uri=CELEX\%3A32024R1689} [Último acceso: 13/09/2026]}. Al tratarse de un proyecto de investigación científica y desarrollo académico, la herramienta queda fuera del ámbito de las obligaciones aplicables a los sistemas de alto riesgo, en virtud de la exclusión que el propio Reglamento prevé para los sistemas desarrollados y puestos en servicio con el único fin de la investigación científica (artículo 2, apartado 6), tal y como se detalla en la evaluación de cumplimiento adjunta en el Anexo~\ref{appendix_ai_act}. No obstante, cabe destacar que el diseño metodológico de este trabajo, basado en el enfoque \textit{human-in-the-loop}, se alinea de forma proactiva con los principios de la normativa europea: exigencia de supervisión humana, transparencia explicativa y mitigación de sesgos discriminatorios.

En lo relativo a la obtención del corpus de estudio, este trabajo se nutre de la recopilación de artículos periodísticos realizada en el marco del proyecto IRIS\_IAMEDIA e InfoIA. Toda la extracción previa de estos datos se llevó a cabo respetando rigurosamente los términos y condiciones de servicio de las cabeceras digitales analizadas. Al tratarse de una investigación académica impulsada por la Universidad Carlos III de Madrid, el procesamiento original de los textos se amparó en las excepciones relativas a la minería de textos y datos (\textit{Text and Data Mining}, TDM por sus siglas en inglés) para fines de investigación científica, heredando este Trabajo de Fin de Máster dichas garantías legales y éticas.

Por último, el desarrollo del código de la investigación se ha realizado en Python con librerías de código abierto. En cuanto a los modelos, se han combinado servicios comerciales accesibles por API con un modelo local de pesos abiertos (\texttt{gemma4:e4b}). La publicación del código y la documentación de la metodología buscan garantizar la transparencia y la reproducibilidad del trabajo.

\newpage